\section{Introduction}
\label{sec:intro}

A looped (or recurrent-depth) language model applies a shared block of layers several times per
token instead of stacking distinct layers
\citep{dehghani2019universal,geiping2025recurrent,zhu2025ouro}. The loop count becomes an
inference-time dial: the same checkpoint can run shallow and cheap or deep and more accurate.
Looped models are now released by several groups, including Ouro \citep{zhu2025ouro}, Dory
\citep{nvidia2026dory} and Nanbeige4.2 \citep{nanbeige2026}, and looping has been shown to help
compositional tool calling specifically \citep{popescu2026toolcalling}. How a deployed model
should set the dial is much less settled.

Two practices dominate. Most deployments fix one loop count for all traffic; Dory's authors, for
example, recommend two loops as a default and four when accuracy matters more
\citep{nvidia2026dory}. The alternative is per-token halting, in the line of adaptive computation
time \citep{graves2016adaptive}, Universal Transformers \citep{dehghani2019universal} and
PonderNet \citep{banino2021pondernet}: Ouro trains an exit gate that can stop each token early
\citep{zhu2025ouro}, and recurrent-depth models can stop when successive loops stop changing the
prediction \citep{geiping2025recurrent}. Per-token exits fit poorly with how models are served.
Requests are decoded in batches, and a batch step costs as much as its deepest token; in our
measurements Ouro's gate is faster than a fixed depth of matching accuracy at batch size 1 but
slower from batch size 4 on, while a depth fixed per request keeps its savings under batching (\S\pending{serving}). Recent
analyses also find that learned halting gates do not track problem difficulty well
\citep{popescu2026halting,olszowiec2026monitoring}.

This paper asks at what granularity, and from what signal, a looped model should choose its
depth, using multi-step tool use as the test bed. Tool use suits the question: a request unfolds
as a sequence of calls, each call is executed, and the final result can be checked exactly, so
every allocation policy can be scored by accuracy and by the compute it spends. We fine-tune
Ouro-2.6B for chained calls, evaluate on the math portion of NESTful \citep{basu2024nestful}
with every intermediate call executed, and compare three granularities of allocation (per tool
call, per request and per token) against fixed depths and against random mixes of fixed depths
at equal compute. Confirmatory tests are fixed in advance and run once on held-out NESTful
requests.

Our findings are:

\begin{itemize}
  \item \textbf{Per-call depth is search, not allocation.} An oracle that picks the best loop
  count for each call reaches .308 on NESTful against .225 for four loops (with an earlier
fine-tuned checkpoint), which looks like a
  large gain. But an oracle that spends the same budget on sampled continuations at four loops
  reaches .325. The same holds for Dory-1.0-4B (.633 vs .617). Varying depth per call mostly
  perturbs the trajectory; it does not allocate compute where it is needed.
  \item \textbf{The useful decision is two loops or four.} On 240 NESTful requests, one loop
  solves .171, two .383, three .442 and four .492 of requests. Three loops is exactly as accurate
  as randomly mixing two- and four-loop runs at the same compute (.441), and one loop is too
  shallow even for single calls. Per request, 85 are solved at both two and four loops, 33 only
  at four and 7 only at two; a router that knew these outcomes would reach .521 at 0.83 times the
  compute of four loops.
  \item \textbf{Depth need is predictable from early loops, not from the request.} Whether a
  request is solved at a given depth is largely stable across sampled runs, yet features of the
  request text predict it no better than chance. Signals from the model itself do, once they are
  read at the right tokens: early loops agree with the final loop on the boilerplate that opens
  a tool call and disagree on the function name and argument values.
  \item \textbf{A deployable escalation router.} Every request starts at two loops; after the
  first ${\sim}45$ tokens, loop-1 and loop-2 signals on those decision tokens decide whether to
  switch to four loops. Nothing extra is computed to decide. \pending{Final held-out result:
  accuracy and compute relative to four loops on NESTful requests 240--719; gain over the best
  random mix with 90\% interval.}
  \item \textbf{A data-leakage audit.} Most MathQA validation problems (1,001 of 1,293 after our
  task construction) have a near-duplicate or an identical computation in the training split.
  Scores on them overstate held-out accuracy (at four loops, .489 on near-duplicates vs .346 on
  clean problems in one evaluation slice), which matters for any work that trains on MathQA and evaluates on
  NESTful, whose math portion comes from the MathQA test split.
\end{itemize}
